\section{Related Work}

Information-gathering problems are often modeled as sequential decision-making problems~\citep{yao2020path}. When there are multiple agents, decentralized information gathering can be viewed as a decentralized POMDP~\citep{bernstein2002complexity}. The dominant approach to Dec-POMDP is to first solve the centralized, offline planning over the joint multi-agent policy space, and then push these policies to agents to execute them in a decentralized fashion~\citep{oliehoek2016concise}. When the state of the environment or agents is not known ahead of time, these approaches become infeasible. Fully decentralized Dec-POMDP solvers exist~\citep{spaan2006decentralized}. However, they require significant memory and incur high computational complexity due to the requirements to compute and store all the reachable joint state estimations~\citep{lauri2020multi}.

Recently, simultaneous distributed approaches based on MCTS have gained significant interest due to their flexibility in trading off computation time for accuracy. The key idea is to use the upper confidence bound (UCB) \citep{kocsis2006improved} for planning the best course of action. To implement cooperation between agents, these methods usually keep a predefined model of the teammates, which can be heuristic or machine learning trained \citep{claes2017decentralised,czechowski2020decentralized,choudhury2021scalable}. However, as they are based on trained knowledge, they are unsuitable for online planning in settings that change unpredictably, such as in disastrous environments.
 
To address this, new approaches based not on apriori knowledge about agents' behavior, but on information sharing among them, have been proposed (Dec-MCTS~\citep{best2019dec}). A key aspect of Dec-MCTS and all its subsequent variations \citep{li2019integrating,li2021dec,nguyen2022multi} is that each agent is assigned a local utility function, which does not measure the total team reward but the contribution of that agent only. To deal with any uncertainty that arises during the mission, Dec-MCTS algorithms allow for online replanning during execution, by having agents update their beliefs about the system. 

Under high uncertainty scenarios, a growing body of literature review in the area of multi-drone systems~\citep{williams2014multi, lizzio2022review, dinelli2023configurations, goeckner2023attrition} explores the significant challenges posed by agent attrition -- the loss or removal of individual drones (due to mechanical failures, environmental factors, and human errors), and highlights the need to address attrition for robust system performance. In systems with attrition in which agents may fail abruptly, all of the above-mentioned approaches do not apply, as they do not allow adjusting to attrition in agents' populations. In the present work, we show that the suboptimality of current Dec-MCTS algorithms is due to the usage of the marginal contribution combined with the submodular properties of the global utility functions.

Another body of literature on related works concerns Open Agent Systems (OASYS), in which agents can enter and leave over time. Most solutions to OASYS are either fully or partially offline, i.e., offline planning with online execution \citep{cohen2017open, chandrasekaran2016individual} or online planning with precomputed offline policies \citep{eck2020scalable}. This is not feasible in applications with significant sources of uncertainty, particularly when the environment's state or mental models of the agents are unknown in advance, and when the agents' failures occur abruptly during execution. A recent work \citep{kakarlapudi2022decision} proposed a fully online approach that leverages communication between agents related to their presence to predict the actions of others. This approach assumes that agents can communicate their existence implicitly. In scenarios where the communication is intermittent or agents vanish without warning, the unforeseeable failure can significantly impact coordination and planning, jeopardizing the overall performance. In addition, the computational complexity of modeling each other's presence and predicting their actions can make the system computationally intractable on a large scale. Thus, it isn't directly equipped to handle sudden agent failures and may require further refinement to be viable in large-scale or highly dynamic environments.

Our paper focuses on problems where agents experience hard failures in an abrupt and unforeseeable fashion. This unpredictable nature makes the existing techniques not directly applicable. Therefore, more research is needed to enhance the robustness and resilience of multi-agent systems in such uncertain attrition scenarios. Our work explicitly tackles this challenge by providing a new approach for online decentralized planning for multi-agents that does not require precomputed offline policies and mental models. Instead, agents reason about the actions and existence of others using directly communicated information.
We employ a computational-effective game-based technique to coordinate agents, enabling adaptive decision-making in the presence of peer failures while ensuring fast convergence in polynomial time relative to system size.